% Fourteen representative nodes; one physical template for all three timelines.
\begin{tikzpicture}[x=1mm,y=1mm,font=\figurefont\fontsize{8}{10}\selectfont,
 history card/.style={draw=FigureBlue!65,fill=white,rounded corners=2mm,line width=.65pt,minimum width=68mm,minimum height=18mm},
 history mark/.style={circle,minimum size=4mm,inner sep=0pt,draw=FigureBlue,fill=FigureBlueFill,line width=.7pt}]
\path[use as bounding box] (0,-5) rectangle (156,184);
\node[anchor=west,font=\figurefont\fontsize{8}{10}\selectfont,text=FigureInk] at (1,181) {神经网络：单元、学习与信息使用};
\node[history card] at (34,169.0) {};
\node[anchor=west,font=\figurefont\fontsize{8}{10}\selectfont,text=FigureBlue] at (4,173.5) {1943\quad 阈值神经元};
\node[anchor=north west,align=left,text width=62mm,font=\figurefont\fontsize{7.5}{9.5}\selectfont,text=FigureInk] at (4,170.5) {简单单元实现逻辑响应\\连接结构组织计算};
\node[history card] at (122,169.0) {};
\node[anchor=west,font=\figurefont\fontsize{8}{10}\selectfont,text=FigureBlue] at (92,173.5) {1958\quad 感知机};
\node[anchor=north west,align=left,text width=62mm,font=\figurefont\fontsize{7.5}{9.5}\selectfont,text=FigureInk] at (92,170.5) {根据分类错误更新权重\\从样本学习识别规则};
\node[history card] at (34,147.5) {};
\node[anchor=west,font=\figurefont\fontsize{8}{10}\selectfont,text=FigureBlue] at (4,152.0) {1980\quad 新认知机};
\node[anchor=north west,align=left,text width=62mm,font=\figurefont\fontsize{7.5}{9.5}\selectfont,text=FigureInk] at (4,149.0) {局部连接形成层次表示\\探索位置变化下的识别};
\node[history card] at (122,147.5) {};
\node[anchor=west,font=\figurefont\fontsize{8}{10}\selectfont,text=FigureBlue] at (92,152.0) {1986\quad 反向传播};
\node[anchor=north west,align=left,text width=62mm,font=\figurefont\fontsize{7.5}{9.5}\selectfont,text=FigureInk] at (92,149.0) {链式法则分配参数梯度\\隐藏层参与整体学习};
\node[history card] at (34,126.0) {};
\node[anchor=west,font=\figurefont\fontsize{8}{10}\selectfont,text=FigureBlue] at (4,130.5) {1997\quad LSTM};
\node[anchor=north west,align=left,text width=62mm,font=\figurefont\fontsize{7.5}{9.5}\selectfont,text=FigureInk] at (4,127.5) {门控调节写入与读取\\记忆路径保留序列信息};
\node[history card] at (122,126.0) {};
\node[anchor=west,font=\figurefont\fontsize{8}{10}\selectfont,text=FigureBlue] at (92,130.5) {1998\quad LeNet};
\node[anchor=north west,align=left,text width=62mm,font=\figurefont\fontsize{7.5}{9.5}\selectfont,text=FigureInk] at (92,127.5) {局部卷积与参数共享\\端到端学习文档识别};
\node[history card] at (34,104.5) {};
\node[anchor=west,font=\figurefont\fontsize{8}{10}\selectfont,text=FigureBlue] at (4,109.0) {2006\quad 深度信念网络};
\node[anchor=north west,align=left,text width=62mm,font=\figurefont\fontsize{7.5}{9.5}\selectfont,text=FigureInk] at (4,106.0) {逐层预训练形成表示\\探索深层网络的初始化};
\node[history card] at (122,104.5) {};
\node[anchor=west,font=\figurefont\fontsize{8}{10}\selectfont,text=FigureBlue] at (92,109.0) {2012\quad AlexNet};
\node[anchor=north west,align=left,text width=62mm,font=\figurefont\fontsize{7.5}{9.5}\selectfont,text=FigureInk] at (92,106.0) {数据与GPU支持深层训练\\卷积表示用于大规模视觉};
\node[history card] at (34,83.0) {};
\node[anchor=west,font=\figurefont\fontsize{8}{10}\selectfont,text=FigureBlue] at (4,87.5) {2014—2015\quad 软对齐注意力};
\node[anchor=north west,align=left,text width=62mm,font=\figurefont\fontsize{7.5}{9.5}\selectfont,text=FigureInk] at (4,84.5) {按当前需求读取源位置\\缓解固定向量的信息瓶颈};
\node[history card] at (122,83.0) {};
\node[anchor=west,font=\figurefont\fontsize{8}{10}\selectfont,text=FigureBlue] at (92,87.5) {2015—2016\quad 残差网络};
\node[anchor=north west,align=left,text width=62mm,font=\figurefont\fontsize{7.5}{9.5}\selectfont,text=FigureInk] at (92,84.5) {捷径传递表示与梯度\\残差分支学习增量变换};
\node[history card] at (34,61.5) {};
\node[anchor=west,font=\figurefont\fontsize{8}{10}\selectfont,text=FigureBlue] at (4,66.0) {2017\quad Transformer／稀疏MoE};
\node[anchor=north west,align=left,text width=62mm,font=\figurefont\fontsize{7.5}{9.5}\selectfont,text=FigureInk] at (4,63.0) {注意力建立跨位置交互\\路由启用部分专家};
\node[history card] at (122,61.5) {};
\node[anchor=west,font=\figurefont\fontsize{8}{10}\selectfont,text=FigureBlue] at (92,66.0) {2018—2020\quad BERT／GPT-3};
\node[anchor=north west,align=left,text width=62mm,font=\figurefont\fontsize{7.5}{9.5}\selectfont,text=FigureInk] at (92,63.0) {预训练学习可复用表示\\双向编码与自回归生成};
\node[history card] at (34,40.0) {};
\node[anchor=west,font=\figurefont\fontsize{8}{10}\selectfont,text=FigureBlue] at (4,44.5) {2021—2022\quad CLIP／S4};
\node[anchor=north west,align=left,text width=62mm,font=\figurefont\fontsize{7.5}{9.5}\selectfont,text=FigureInk] at (4,41.5) {图文配对学习共同表示\\结构化状态处理长序列};
\node[history card] at (122,40.0) {};
\node[anchor=west,font=\figurefont\fontsize{8}{10}\selectfont,text=FigureBlue] at (92,44.5) {2023—2024\quad Mamba／Mamba-2};
\node[anchor=north west,align=left,text width=62mm,font=\figurefont\fontsize{7.5}{9.5}\selectfont,text=FigureInk] at (92,41.5) {输入选择影响状态更新\\联系状态空间与注意力};
\begin{scope}[on background layer]
\draw[FigureMuted!35,line width=2.2mm,line cap=round] (73,169.0) -- (83,169.0) .. controls (91,164.0) and (65,152.5) .. (73,147.5) -- (83,147.5) .. controls (91,142.5) and (65,131.0) .. (73,126.0) -- (83,126.0) .. controls (91,121.0) and (65,109.5) .. (73,104.5) -- (83,104.5) .. controls (91,99.5) and (65,88.0) .. (73,83.0) -- (83,83.0) .. controls (91,78.0) and (65,66.5) .. (73,61.5) -- (83,61.5) .. controls (91,56.5) and (65,45.0) .. (73,40.0) -- (83,40.0);
\draw[white,dashed,line width=.3mm,dash pattern=on 1.5mm off 1.3mm] (73,169.0) -- (83,169.0) .. controls (91,164.0) and (65,152.5) .. (73,147.5) -- (83,147.5) .. controls (91,142.5) and (65,131.0) .. (73,126.0) -- (83,126.0) .. controls (91,121.0) and (65,109.5) .. (73,104.5) -- (83,104.5) .. controls (91,99.5) and (65,88.0) .. (73,83.0) -- (83,83.0) .. controls (91,78.0) and (65,66.5) .. (73,61.5) -- (83,61.5) .. controls (91,56.5) and (65,45.0) .. (73,40.0) -- (83,40.0);
\end{scope}
\node[history mark] at (73,169.0) {};
\draw[-{Stealth[length=1.4mm]},FigureStroke,line width=.6pt] (75.5,169.0) -- (80.5,169.0);
\node[history mark] at (83,169.0) {};
\node[history mark] at (73,147.5) {};
\draw[-{Stealth[length=1.4mm]},FigureStroke,line width=.6pt] (75.5,147.5) -- (80.5,147.5);
\node[history mark] at (83,147.5) {};
\node[history mark] at (73,126.0) {};
\draw[-{Stealth[length=1.4mm]},FigureStroke,line width=.6pt] (75.5,126.0) -- (80.5,126.0);
\node[history mark] at (83,126.0) {};
\node[history mark] at (73,104.5) {};
\draw[-{Stealth[length=1.4mm]},FigureStroke,line width=.6pt] (75.5,104.5) -- (80.5,104.5);
\node[history mark] at (83,104.5) {};
\node[history mark] at (73,83.0) {};
\draw[-{Stealth[length=1.4mm]},FigureStroke,line width=.6pt] (75.5,83.0) -- (80.5,83.0);
\node[history mark] at (83,83.0) {};
\node[history mark] at (73,61.5) {};
\draw[-{Stealth[length=1.4mm]},FigureStroke,line width=.6pt] (75.5,61.5) -- (80.5,61.5);
\node[history mark] at (83,61.5) {};
\node[history mark] at (73,40.0) {};
\draw[-{Stealth[length=1.4mm]},FigureStroke,line width=.6pt] (75.5,40.0) -- (80.5,40.0);
\node[history mark] at (83,40.0) {};
\node[anchor=west,font=\figurefont\fontsize{8}{10}\selectfont] at (1,27) {持续研究的问题};
\node[draw=FigureBlue,fill=FigureBlue!8!white,rounded corners=2mm,minimum width=49mm,minimum height=23mm,line width=.65pt] at (25,10) {};
\node[font=\figurefont\fontsize{8}{10}\selectfont] at (25,16) {能力与数据};
\node[align=center,font=\figurefont\fontsize{7.5}{9.5}\selectfont] at (25,8) {自监督、多模态\\数据覆盖与迁移};
\node[draw=FigureGreen,fill=FigureGreen!8!white,rounded corners=2mm,minimum width=49mm,minimum height=23mm,line width=.65pt] at (78,10) {};
\node[font=\figurefont\fontsize{8}{10}\selectfont] at (78,16) {结构与成本};
\node[align=center,font=\figurefont\fontsize{7.5}{9.5}\selectfont] at (78,8) {长序列、稀疏专家\\状态存储与通信};
\node[draw=FigureYellow,fill=FigureYellow!8!white,rounded corners=2mm,minimum width=49mm,minimum height=23mm,line width=.65pt] at (131,10) {};
\node[font=\figurefont\fontsize{8}{10}\selectfont] at (131,16) {训练与可靠性};
\node[align=center,font=\figurefont\fontsize{7.5}{9.5}\selectfont] at (131,8) {后训练、推理时计算\\校准与稳健性};
\end{tikzpicture}
